\documentclass[11pt,letterpaper]{article}
\usepackage{amssymb}
\usepackage[margin=.5in]{geometry}

\setlength{\parindent}{0pt}
\begin{document}

%\begin{center}
%{\Large \bf Winter 24-25 :: CR 04 :: Project }\\
%{due Tu. 01/21/25, by email, before class }
%\end{center}

\small

{\bf Rubric for SEP}

Pour SEP, on utilise le bon algo de LAPACK. Celui là divise les FLOPS par 2 et tiens en compte la symmétrie, 
c'est pas évident à faire sinon en faisant un "Hessenberg symmétrique". Donc il est plus facile de faire 
un "Hessenberg symmétrique" mais c'est pas bon.

Shift are computed by hand.

cleanup: For NEP, we reduced to standard form with "unitary" transformation.
cleanup: For SVD, make sure the singlar values are nonnegative!!!!

A good way to show the cubic/quadratic convergence and the deflation window that moves is simply to 
print a small matrix for all the iterations.

Les matrices orthogonales sont bien calculés et la representativité est conservé tout du long

Il n'y a pas de extra n^3 en FLOPS pas de extra n^2 en storage, par exemple:
        U1,B1,V1 = biDiag(M,n,m,f)
        U2,D,V2 = biDiagToDiag(B1,it,epsilon,n,m,f)
        U = U1 * U2
Non!!!!



le test est fait sur une matric D ou T qui est avec des zeros, ce n'est pas la matrice qui vient de itérations!!!!
Les 1e-16 de la sous diagonal doivent être mis explicitement à zeros, et pas que les 1e-16, ...

The construction of Q is done using a "reverse" algorithm and is not done on teh fly. That way we use less flops.



( part 1 ) [ + 5 points ] \\
A basic code is working. It is based on a reduction to condensed form, either
Hessenberg form or symmetric tridiagonal form. Once a symmetric tridiagonal
form is computed then a QR iteration algorithm (possibly without shift),
possibly with $\mathcal{O}(n^3)$ iteration per step (as opposed to
$\mathcal{O}(n)$ iteration per step) is performed until convergence to diagonal
form. The checks work and are based on $\| A – Q D Q^T \|_1 / \| A \|_1$, and 
$\| I – QQ^T \|$ and the fact that $D$ is a diagonal matrix. The code work on random
symmetric small size matrices (e.g. $n = 20$).

( part 2 ) [ + 2 points ] \\
A shift is added. The shift enables cubic convergence on the eigenvalues in
position $t_{n,n}$. The shift is computed using ``Wilkinson shift''. The
iterations now look like $Q, R  = \textmd{qr} ( T – \mu I )$, and $T = RQ + \mu * I$. 

( part 3 ) [ + 4 points ] Deflation is done, that is that when an eigenvalue is
converged, then we need to work on a smaller window. Also, it is important to
scan the subdiagonal for possible zeros, so that there is no ``near zero''
numbers on the subdiagonal of the active matrix. The criteria to deflate or
decide on reducibility is 
$$| t_{i-1,i} | < \varepsilon ( | t_{i-1,i-1} | + | t_{i,i} | ).$$
There is two parts: (1) one reducing from the bottom right, and (2) reducing from the top left.
But then also come back to the unconverged block on the top left.

( part 4 ) [ this is optional, and can bring + 3 points ] \\
The two steps QR, RQ on symmetric tridiagonal matrices are optimized to only do
$\mathcal{O}(n)$ iterations per step. (And $\mathcal{O}(n^2)$  if eigenvectors are
required.) (As opposed to $\mathcal{O}(n^3)$.)

( part 5 ) [ + 5 points, this is the good way to do it and if done correctly
replace part 4, so you cannot have point for part 4 and part 5 at the same time
] \\
The QR iteration is performed using the implicit Q theorem and bulge chasing.

( part 6 ) [ +1 point ] \\
A flag is present to compute either ( eigenvalue only ) or ( eigenvalue and eigenvector )

( part 7 ) [ + 3 points ] \\
Experimental part. Some experiments are performed: (*) number of iterations per
eigenvalues, (*) 

( Holistic criteria ) [ +2 points ] \\
software engineering is well done, (good interfaces, good data structure, etc.)
code is readable, interview/one-on-one meeting was well done, student was able
to answer questions and present their code


the checks are done with respect to a cleaned up T or D, not the matrix A on which we iterate!




Choose one of the three problems below. Do not do all three. You will only be
graded on one problem. The one you want me to grade.

Please ask me by email if directions below are not clear.

Codes need to be in ``julia''.
 
\section*{Logistic}

Trois sujets au choix: NEP, nonsymmetric eigenvalue problem: matrice $n \times
n$, calculer ``Schur decomposition''.  SEP, symmetric eigenvalue problem:
matrice $n \times n$, calculer $Q D Q^T$.  SVD, singular value decomposition:
matrix $m \times n$, calculer $U S V ^T$.  Prenez celui que vous préférez, plus
d’indications sur les sujets en dessous. Pour SVD, le livre de Darve et Wooters
est très bien fait (Section 5.4). Normalement vous suivez le livre et c’est
bon. Pour le NEP, il y a déjà beaucoup de fait dans le ./src/gees.jl sur le
GitHub de Eric Darve, évidement, il ne s’agit pas de simplement copié-collé les
codes et me donner cela. Oui vous pouvez tout réutilisé. Et je vous encourage à
bien tout regarder, mais il y a du boulot en plus à faire. Pour SEP, il faut
"symmétriser NEP".  Pour SEP et SVD, les problèmes sont plus simples à la base,
mais vous commencez à zero. Ce seraient mieux que, sur les 11 étudiants, on ait
un peu de trois sujets, mais choisissez ce que vous plait.
 
Les projets doivent être fait individuellement. Vous pouvez communiquer entre
vous, mais pas de recopiage de code. Je n’ai pas eu de problème de triche pour
les DM1, DM2 et DM3. Donc je ne suis pas inquiet pour les projet de fin de
semestre. Continuez comme cela pour le travail.
 
Dû: Le rapport écrit est simplement vos codes. Soit des .jl ou un .ipynb. (Ou
un mix.) Tout est dû le mardi 21 janvier au matin avant le cour. Vous m’envoyez
tout celà par email. Attention: le DM5 sera aussi dû ce jour là.  Les
présentations seront de 15 minutes. Et vous expliquerez vos codes à la classe. (Moi
inclus.) Des questions seront posées. Le mieux est de ne pas trop stresser pour
la presentation, si vous avez écrit vos codes et les comprenez cela se passera
bien.  Ordre de passage pour les presentations: a priori: tirage au sort mardi
21 janvier au matin.

La note sera une note ``holistique'' donc ce n’est 10\% pour l’oral, 40\% pour
l’écrit, mais c’est une note pour le tout. Principalement au final ce qui
compte ce sont les codes, mais si vous ne pouvez pas les expliquer clairement
pendant l’oral, cela ne va pas aller.
 
Le projet de fin de semestre est 50\% de votre note finale. (L’autre 50\% est
fait des 5 DMS, 10\% chacun.)
 
Si vous êtes bloqués ou avez des questions pendant les vacances, n’hésitez pas
à m’envoyer un email. Je peux voir si j’ai le temps de répondre. Je pense que
si vous ne travaillez pas pendant les vacances sur les projets, évidemment ce
sera plus chargé à la rentrée mais c’est jouable.
 
\section*{Project \#1: NEP, nonsymmetric eigenvalue problem}

In input a nonsymmetric $n$-by-$n$ matrix $A$. In output the Schur form of $A$.
( $A = Q T Q^T$. ) Only real arithmetic, no complex arithmetic. $T$ is quasi
triangular so 1x1 block for real eigenvalues and 2x2 block for complex
eigenvalues. Please have a look at gees.jl in ./src of GitHub of Eric Darve for
a template. Algorithms needs to be the implicit Q QR iteration algorithm using
Francis shift. Please reduce to Hessenberg form initially, then each iteration
needs to be $\mathcal{O}(n^2)$ and so needs to take into account the structure of H.
Three things not done in gees.jl that you need to do (1/3) Need flags for
either 'D' compute the diagonal only, 'T' compute the Schur form only, or 'V'
compute both $Q$ and $T$. (2/3) Need to fix the ``explicit shift strategy''.
(3/3) Return the 2x2 block (whenever they exist on the diagonal) as ``standard''
form. See DM3 question 3.

Goal is that your code is able to work on a random matrix. Then you can try on
matrix $A = VDV^{-1}$ when $V$ is ill conditioned. Also try with cluster of
eigenvalues. A few other types of matrices are interesting.

Analysis: (*) observe quadratic convergence for bottom right side of matrices,
(*) return a metric such as ``number of QR iterations per eigenvalue'', (*) time
your code for various $n$ and try to scale up to larger $n$.

Rubrique: your code works for random matrices, is clearly written
and the analysis (stability, convergence and timing) is done, the choice made
in the design of the code are clearly explained in the comments of the code, you are able to clearly 
explain the code and answer questions 
during the oral presentation.

% How to get more points:

% (*) From the Schur form, write the computation for eigenvectors. Check that you
% have $A = V D V^{-1}$ (if $A$ is diagonalizable ``enough'').

% (*) Reorder the eigenvalues in the Schur form (and that associated Schur
% vector) such that all eigenvalues with positive real parts are at the start of
% the Schur form, and all the eigenvalues with negative real parts are at the
% end. 

% (*) Explore EAD: Early Aggressive Deflation

\section*{Project \#2: SEP, symmetric eigenvalue problem}

In input a symmetric $n$-by-$n$ matrix $A$. ($A=A^T$.) In output the
eigendecomposition of $A$. ( $A = Q D Q^T$. ) $D$ is diagonal real and $Q$ is
$n$-by-$n$ orthogonal: $Q^TQ = QQ^T = I$. Only real arithmetic, no complex
arithmetic. (So do not take $A$ as complex Hermitian ($A=A^H$). Real symmetric
is enough. First step is to take A and reduce to symmetric tridiagonal form.
You can only use the lower part of $A$. (Or the upper, as you wish. But only
one half. I recommend starting with gehrd! and then adapting for the symmetric
case. But as you wish.) The nonsymmetric QR algorithm that is explained in Darve and Wootters 
needs to be ``symmetrized''. 

Need flags for
either ('D') compute the diagonal only or 'V'
compute both $Q$ and $D$. 
The goal is to perform $\mathcal{O}(n)$ operations per QR step if eigenvalues only are requested,
and $\mathcal{O}(n^2)$ operations per QR step if eigenvalues and eigenvectors are requested.

Analysis: (*) observe cubic convergence for bottom right side of matrices,
(*) return a metric such as ``number of QR iterations per eigenvalue'', (*) time
your code for various $n$ and try to scale up to larger $n$.

Rubrique: your code works for random matrices, is clearly written
and the analysis (stability, convergence and timing) is done, the choice made
in the design of the code are clearly explained in the comments of the code, you are able to clearly 
explain the code and answer questions 
during the oral presentation.

% How to get more points:

% (*) Write the code for complex Hermitian matrix $A$ in input. The code needs to reduce
% $A$ (which is complex Hermitian) to a real symmetric tridiagonal matrix.
% Then you can reuse your previous code for a real symmetric tridiagonal matrices.

% (*) Explore EAD: Early Aggressive Deflation

% (*) Explore Divide and Conquer method.

\section*{Project \#3: SVD, singular value decomposition}

In input an $m$-by-$n$ matrix $A$. (Please take $m\geq n$.) In output the thin SVD
of $A$. ( $A = U S V^T$. ) Only real arithmetic, no complex arithmetic. Please
read section 5.4 in Darve and Wooters. The algorithm is explained there.

Analysis: (*) observe cubic convergence for bottom right side of matrices,
(*) return a metric such as ``number of QR iterations per singular value'', (*) time
your code in the square case for various $m=n$ and try to scale up to larger $m=n$.

Rubrique: your code works for random matrices, is clearly written and the
analysis (stability, convergence and timing) is done, the choice made in the
design of the code are clearly explained in the comments of the code, you are able to clearly 
explain the code and answer questions 
during the oral presentation.

% The code needs to work for $m>n$

% How to get more points:

% (*) Write the code for complex matrix $A$ in input. The code needs to reduce
% to $A$ (which is complex ) to a real symmetric tridiagonal matrix.
% Then you can reuse your previous code for areal symmetric tridiagonal matrices.

% (*) Explore EAD: Early Aggressive Deflation

\end{document}
